\documentclass[11pt,a4paper]{article}
\usepackage[T1]{fontenc}
\usepackage[utf8]{inputenc}
\usepackage{lmodern}
\usepackage[margin=22mm,headheight=15pt]{geometry}
\usepackage{microtype}
\usepackage{amsmath,amssymb,amsthm,mathtools}
\usepackage{booktabs,tabularx,longtable,array}
\usepackage{xcolor}
\usepackage{enumitem}
\usepackage{xurl}
\usepackage{hyperref}
\usepackage{fancyhdr}
\usepackage{needspace}
\definecolor{ink}{HTML}{17334B}
\definecolor{accent}{HTML}{176C72}
\hypersetup{colorlinks=true,linkcolor=ink,citecolor=accent,urlcolor=accent,
 pdftitle={Beyond the Dirichlet GRH endpoint: methods and further targets},
 pdfauthor={Research analysis prepared for Vladimir Reshetnikov},
 pdfsubject={Audit and transfer analysis of grh-lean with complete derived proofs and certificates},
 bookmarksnumbered=true,hypertexnames=false}
\pagestyle{fancy}
\fancyhf{}
\fancyhead[L]{\small Beyond the Dirichlet GRH endpoint}
\fancyhead[R]{\small 25 November 2026}
\fancyfoot[C]{\thepage}
\setlength{\parskip}{4pt}
\setlength{\parindent}{0pt}
\setlength{\emergencystretch}{2em}
\setlist{itemsep=3pt,topsep=5pt,leftmargin=*}
\newcommand{\Q}{\mathbb Q}
\newcommand{\Z}{\mathbb Z}
\newcommand{\R}{\mathbb R}
\newcommand{\C}{\mathbb C}
\newcommand{\N}{\mathbb N}
\newcommand{\file}[1]{\nolinkurl{#1}}
\newcommand{\repo}{\texttt{grh-lean}}
\newtheorem{theorem}{Theorem}[section]
\newtheorem{lemma}[theorem]{Lemma}
\newtheorem{proposition}[theorem]{Proposition}
\newtheorem{corollary}[theorem]{Corollary}
\theoremstyle{definition}
\newtheorem{definition}[theorem]{Definition}
\theoremstyle{remark}
\newtheorem{remark}[theorem]{Remark}
\newcolumntype{Y}{>{\raggedright\arraybackslash}X}
\numberwithin{equation}{section}
\begin{document}
\hypersetup{pageanchor=false}
\begin{titlepage}
{\color{accent}\large RESEARCH REPORT \hfill 25 November 2026\par}
\vspace{20mm}
{\color{ink}\Huge\bfseries Beyond the Dirichlet\\[3pt] GRH endpoint\par}
\vspace{8mm}
{\LARGE Transferable arithmetic, a nonabelian certificate,\\[4pt]
and additional class-group exclusions\par}
\vspace{12mm}
{\large Analysis of \href{https://github.com/lsernau/grh-lean}{\texttt{lsernau/grh-lean}}\\
and five preceding reports in \texttt{ProveIt/docs/incoming}\par}
\vspace{8mm}
Prepared for Vladimir Reshetnikov

\vspace{9mm}
\begin{abstract}
The repository proves RH and GRH for Dirichlet \(L\)-functions. Its internal
arithmetic is more reusable than that endpoint alone suggests. We extract
a uniform torsion-selection theorem and extend it, with a proof, to arbitrary
families of bounded power series over \(\mathcal O_{\C_p}\) with finite total
valuation budget. A single finite torsion grid controls every determinant
simultaneously, independently of degrees and coefficient fields.

For a genuinely nonabelian target, we construct an explicit positive
\(A_5\) character detector. A fully normalized one-sided prime-sum estimate
would imply RH for a specified degree-60 splitting field and three subfields;
the required arithmetic estimate remains unproved. We quantify why fixed
abelian congruence weights cannot supply these coefficients.

Finally, an intermediate genus-field argument answers an unfinished question
in a prior report and gives four additional exclusion families for imaginary
quadratic class groups of exponent six. These are partial classification
results, not a complete census theorem. Full mathematical proofs, exact
finite certificates, pinned source evidence, and explicit limitations
accompany the analysis. No new Lean verification or worldwide novelty claim
is made.
\end{abstract}
\vfill
{\small\color{ink}
Source revision: \file{72fead102f663d5c7a407964a13dc0cf9d977207}\\
Prior-report repository revision: \file{a92443f00663b1b8b3c233d17df0ad49dc4f2b2c}\\
Assessment date uses America/Los\_Angeles. All external sources were consulted
during this analysis.}
\end{titlepage}
\hypersetup{pageanchor=true}

\section*{Results at a glance}
\addcontentsline{toc}{section}{Results at a glance}
The ideas do yield a proved reusable extension and new partial exclusions
beyond the usual consequences of RH. The present analysis does not prove
a further major named conjecture in full. The distinction between a
completed deduction and a missing arithmetic input is essential to
interpreting the results.

\begin{center}
\small
\begin{tabularx}{\textwidth}{@{}>{\raggedright\arraybackslash}p{31mm}YY@{}}
\toprule
Direction & Completed here & Remaining requirement\\
\midrule
Bounded \(p\)-adic analytic families &
A single bounded-order torsion tuple controls total positive valuation
loss for any family with finite witness budget; matrix and normalized-length
corollaries. &
To settle a Selmer or Iwasawa question, construct the relevant integral
analytic presentation and prove the witness budget.\\
\addlinespace
Nonabelian \(A_5\) RH &
Positive mean-one detector, exact treatment of inertia and prime powers,
one-sided analytic implication, explicit field, and a mean-square
obstruction to fixed congruence substitution. &
Prove the new detector's square-root-scale one-sided prime-sum estimate.\\
\addlinespace
Exponent-six quadratic class groups &
An intermediate genus-field certificate and four stronger necessary
conditions in the remaining census window. &
Enumerate or eliminate every surviving discriminant to prove the full
432-field classification.\\
\addlinespace
Proposed stronger torsion rules &
An exact counterexample to a maximal-order, no-loss principle with
arbitrary ramified coefficients. &
Any successful replacement must retain the appropriate conductor,
ramification, or error hypotheses.\\
\bottomrule
\end{tabularx}
\end{center}

The first direction genuinely transfers an internal ingredient without
using the final RH theorem. The second transfers the analytic architecture
to a problem outside the exported Dirichlet statements. The third is a
new refinement of an application in the earlier reports: it uses the
exported theorem, so it is labeled separately rather than presented as
an extension of the theta construction.

The main mathematical sections are self-contained:
Section~\ref{sec:torsion-transfer} proves the torsion theorem;
Section~\ref{sec:a5-transfer} proves the nonabelian reduction;
Section~\ref{sec:genus-extension} proves the class-group exclusions.
Section~\ref{sec:priorities} identifies the precise next research gates.
The supporting Python programs certify finite computations; they do not
substitute for the proofs of the infinite statements.

\clearpage
\begingroup
\small
\setlength{\parskip}{0pt}
\tableofcontents
\endgroup
\clearpage

\section{What the repository proves, and what this audit establishes}
\label{sec:audit}

\subsection{The endpoint is Dirichlet GRH}

We inspected the pinned source, not just the repository title.
The file \file{grh/src/Statements.lean} imports only Mathlib and states
the propositions proved by \file{grh/src/Main.lean}. The five exports are
listed below. In the general Dirichlet statement, \(N\ne0\), all characters
modulo \(N\) are allowed, and the pole of the trivial character at \(s=1\)
is excluded explicitly. The zero-on-line statement concerns
\(0<\Re s<1\). The separate zeta statement excludes its usual trivial
zeros. These are substantive statements, including imprimitive
characters, with no custom redefinition of GRH hiding in the
endpoint file~\cite{repo}.

\begin{center}
\small
\begin{tabularx}{\textwidth}{@{}p{76mm}Y@{}}
\toprule
Export in \file{Main.lean} & Mathematical scope\\
\midrule
\file{riemannHypothesis} & Mathlib's proposition \file{RiemannHypothesis}.\\
\file{riemannZeta_eq_zero_re} & Every zeta zero apart from the stated exclusions
has real part \(1/2\).\\
\file{LFunction_ne_zero_of_half_lt_re} &
Dirichlet \(L(\chi,s)\ne0\) for \(\Re s>1/2\), apart from the pole case.\\
\file{LFunction_eq_zero_re} &
Dirichlet zeros in the open critical strip lie on the critical line.\\
\file{LFunction_eq_zero_re_of_isPrimitive} &
The corresponding statement restricted to primitive characters.\\
\bottomrule
\end{tabularx}
\end{center}

This immediately supplies RH for abelian Dedekind zeta functions by their
Dirichlet factorization. It is not an export of arbitrary nonabelian
Dedekind, Hecke-over-number-fields, Artin, or automorphic GRH.
That scope is exactly why an \(A_5\) splitting field is a meaningful
further target.

\subsection{Reproducibility and trust boundary}

The readable development uses Lean 4.35.0 and Mathlib revision
\file{56a005855ea17e0b0958f73ec801ae140b0a1679}.
At the pinned commit, our source inventory finds 803 readable-edition
Lean files under \file{grh/src}, totaling 477,124 physical lines including
comments, and 794 compact-edition files totaling 112,827 lines.
The accompanying script records hashes of key files and the declaration
locations used in this report.

The repository and independent checking pages report successful builds,
kernel replay or independent checking, endpoint comparison, and only
\file{propext}, \file{Classical.choice}, and \file{Quot.sound} for the
exported theorems~\cite{repo,sutherland-check,sonoda-check,rubinstein-check}.
This report accepts the mathematical endpoint specified by the task and
uses those records as external verification evidence. It does not claim
to have repeated them. No Lean toolchain was installed in this workspace,
and no full \texttt{lake build}, \texttt{\#print axioms}, or independent
kernel replay was run here.

Our limited lexical scan strips comments and strings and finds no
standalone \file{sorry}, \file{admit}, \file{axiom}, or \file{unsafe}
tokens in either edition's \file{src} tree. That is useful source hygiene,
not a kernel audit or a proof about dependency trust. An expected
placeholder in a challenge file outside those source trees is not evidence
of an unfinished theorem in the compiled endpoint.

The three mathematical developments below were separately cross-checked
for their substantive algebraic and analytic steps. The finite numerical
certificates were executed. The derived extensions are written proofs
in this report; none is represented as a new Lean-checked theorem.

\section{Which ideas can actually be transplanted?}
\label{sec:mechanism}

\subsection{The arithmetic comparison that produces the saving}

The proof's distinctive content is the construction and comparison of
arithmetic lattices of forms. The final Landau argument is an analytic
consumer of an arithmetic estimate. It cannot on its own generate the
estimate for another \(L\)-function.

At a schematic level, short integral forms with a controlled index give
an upper bound for a normalized covolume. A sufficiently large independent
family makes the useful prime remainder visible with a substantial
coefficient. Long-form estimates from theta chains, congruences, and
Selmer-controlled twists give a lower bound with the same leading term.
In the source's notation the comparison has the shape
\[
 M-\text{controlled error}
 \ \le\ \frac{V+2\delta_{\mathrm{lat}}}{d}
 \ \le\ M-\frac nd\,mR+\text{controlled error}.
\]
Here \(M\) is the common main term, \(d\) the target dimension, \(n\) the
spanning dimension, \(V\) a covolume term, and
\(\delta_{\mathrm{lat}}\) an index contribution. The subtraction of the
common \(M\) leaves an upper bound on \(mR\).
This formula is explanatory notation, not a replacement for the exact
Lean inequalities~\cite{repo,alpoge-remark}.

\begin{center}
\small
\begin{tabularx}{\textwidth}{@{}>{\raggedright\arraybackslash}p{28mm}YY@{}}
\toprule
Stage & Source location & Role in a transfer\\
\midrule
Short and enough &
\file{New/Reduction/ShortAndEnough.lean} &
Integral short vectors, index control, and enough independent vectors.\\
\addlinespace
Local-to-global lattice &
\file{New/Reduction/Lattices/TwistedLattice.lean} &
Product formula converts locally chosen twists and denominator bounds
into a global covolume inequality.\\
\addlinespace
Long forms &
\file{New/Determinant/Twist/LongForms.lean} &
Arithmetic lower bound whose main term matches the short-form bound.\\
\addlinespace
Cancellation and scale &
\file{New/Reduction/MainInequality.lean};
\file{New/Reduction/OneSidedBound.lean} &
\file{star_inequality} gives
\(m\,\mathrm{totalRemainder}\le fa(\log b)^5\), then \(m\asymp\sqrt{x}\).\\
\addlinespace
Zeros &
\file{New/Reduction/Zeros/Criterion.lean};
\file{Lib/Analytic/Landau.lean} &
A one-sided bound and appropriate real residue signs exclude poles of
a transform, hence unwanted zeros.\\
\bottomrule
\end{tabularx}
\end{center}
All paths in this table are relative to \file{grh/src}.

In particular, an intermediate bound in
\file{eventually_primeRemainder_natCast_le} is
\[
 R_q(x)\le \bigl(64\varphi(q)+3\bigr)x^{-1/2+6\varepsilon}
 \qquad\text{eventually}.
\]
The choice of \(\varepsilon\), the fixed modulus, the auxiliary scales,
and the source's prime-number-theorem input lead to the final requested
one-sided exponent. The square-root scale comes from the length of
the arithmetic construction; it is not supplied by an abstract
zero-free principle.

The actual prime coefficients involve congruence conditions and a
quadratic factor of the form
\[
 \Lambda(n)\bigl(1-\chi_K(n)\bigr)
             \mathbf 1_{\,n\equiv1\pmod q}.
\]
The source also selects auxiliary quadratic twists to prevent cancellation
among residues in a prescribed box. A replacement target must justify
both its coefficients and its noncancellation input. Relabeling
\(L(\chi,s)\) as a general automorphic \(L\)-function does neither.

\subsection{The local quantifiers matter}

The twisted-lattice theorem permits the local twist to depend on the
rational prime and on the embedding into a \(p\)-adic algebraic closure.
Its product-formula conclusion does not assert that a single global
character works at every finite place. Similarly, selecting factors
before taking their products or gcds can change a divisibility statement.
The earlier reports already emphasize this quantifier issue; it must
survive every attempted transfer.

The torsion theorem in Section~\ref{sec:torsion-transfer} has a different,
particularly useful uniformity: at one fixed prime it chooses a finite
grid before seeing the family of polynomials or analytic functions.
That uniformity makes an infinite-family extension possible. A static
import-closure scan of its source file reaches 103 local modules,
58,528 lines, and 3,388,790 bytes, with Mathlib as the only external
import name. This is the current syntactic import closure, not a
minimal mathematical dependency set. It contains no final RH endpoint
and makes this theorem a plausible standalone extraction target.

\subsection{What the five preceding reports already cover}
\label{sec:prior}

All five ZIPs in the selected \file{docs/incoming} snapshot were inspected.
Their hashes and exact internal report paths are recorded in
\file{evidence/prior_reports_manifest.json}. The table summarizes their
claimed completed arguments and their explicit gaps; enumerations
recorded in those packets were not all rerun in this analysis~\cite{prior}.

\begin{center}
\footnotesize
\begin{tabularx}{\textwidth}{@{}p{49mm}YY@{}}
\toprule
Archive & Already developed there & Boundary relevant here\\
\midrule
\file{grh_transfer_analysis.zip} &
Determinant sandwich, scalar Landau reduction, conductor-23 cubic target,
and seven small algebraic Lean lemmas. &
No principal-class prime bound and no nonabelian RH conclusion.\\
\addlinespace
\file{grh_transfer_analysis_2026-11-25.zip} &
Higher-order pole criterion, sparse witness families, Tate-inertia
obstruction, and product/gcd quantifier analysis. &
No arithmetic construction of the sparse family; no Vandiver or
\(K\)-group vanishing theorem.\\
\addlinespace
\file{grh-transfer.zip} &
A reported 229-field abelian sufficient-condition classification,
small-prime representations, split-prime criteria, and upper limsup. &
Reported enumeration not rerun here; an equivalent prime inequality
is not a proved inequality.\\
\addlinespace
\file{grh-lean-transfers.zip} &
A reported 26-field quintic CM classification, genus containment,
weighted races, and \(S_3\) detector analysis. &
No \(S_3\) zero theorem; algebraic classifications do not extend the
theta construction.\\
\addlinespace
\file{grh_transfers.zip} &
Partial exponent-six census, mod-4 window positivity, eta-product
variance, and a conductor-23 sign equivalence. &
The full exponent-six census remains unproved; it explicitly asks for
intermediate genus-field improvements.\\
\bottomrule
\end{tabularx}
\end{center}

The present report adds the analytic and infinite-family torsion theorem,
the \(A_5\) construction with its exact coefficients and variance, and the
intermediate genus-field improvement. It does not relabel the previous
scalar Landau theorem, the previously observed Tate-inertia obstruction,
or the existence of simultaneous nonvanishing results as new discoveries.

For comparison, combining RH with the established nonnegativity of the
de Bruijn--Newman constant gives \(\Lambda=0\) immediately. That is an
important consequence but not a transfer of the internal method, and it
is excluded from this report's new findings~\cite{rodgers-tao}.
The same restraint applies to familiar effective-prime, class-number,
and conditional algorithm consequences unless an additional argument
substantially improves the preceding analysis.


\clearpage
\section{An arithmetic transfer: one torsion point for many analytic constraints}
\label{sec:torsion-transfer}

The most portable arithmetic ingredient we found is not the final RH
statement. It is a selection theorem buried in
\texttt{New/Congruences/Selmer/RootsOfUnity.lean}. Its root-counting argument
controls the divisibility of many polynomials at the same tuple of roots of
unity, with an order bound independent of their degrees. This section
extracts a slightly stronger budget from that argument and proves its
extension to bounded power series and infinite families. These are derived
results of this analysis; no claim of literature priority is made.

Fix a prime $p$. Write $v$ for the valuation on $\C_p$, normalized by
$v(p)=1$, and set $v(0)=+\infty$. Its ring of integers is
$\mathcal O_{\C_p}=\{z:v(z)\ge0\}$. For a finite real number $t$, define
$(a-t)_+=\max(a-t,0)$, with $(+\infty-t)_+=+\infty$.
For nonnegative numbers indexed by an arbitrary set, a sum means the
supremum of its finite partial sums. In particular, no countability
assumption is implicit in the statements below.

\subsection{The finite polynomial mechanism and its quantitative budget}

The source theorem
\texttt{New.Selmer.exists\_rootsOfUnity\_ellVal\_eval\_le\_add}
selects one tuple for a finite family of integral polynomials. Its
intermediate lemmas retain more information than its final statement.
We give the argument, including that information, so the analytic
extension has a precise mathematical starting point.

For an integral polynomial $g(T)=\sum_k a_kT^k$ and $r>0$, put
\[
 v_r(g)=\min_k\{v(a_k)+kr\}.
\]
This is the Gauss valuation of radius $p^{-r}$. It is additive on
products, and $v(g(y))\ge v_r(g)$ when $v(y)\ge r$.
Set
\[
 r_h=\frac1{\varphi(p^h)}\qquad(h\ge1).
\]
Every primitive $p^h$-th root of unity satisfies $v(\zeta-1)=r_h$.

\begin{lemma}[One coordinate, with an explicit choice of orders]
\label{lem:one-coordinate-torsion}
Let $\rho>0$, $0<\delta<1$, and $C\ge0$. Suppose that $H$ satisfies
\[
 C<\delta\sum_{\substack{1\le h\le H\\r_h\le\rho}}\frac1{h+1}.
\]
If $g\in\overline{\Z}_p[T]$ and $v_\rho(g)\le C$, there are an
admissible $h\le H$ and a primitive $p^h$-th root of unity $\zeta$ such
that
\[
 v(g(\zeta-1))\le v_{r_h}(g)+\delta.
\]
Here $\overline{\Z}_p$ denotes the integral elements of
$\overline{\Q}_p$.
\end{lemma}

\begin{proof}
For each admissible $h$, let $R_h$ be the multiset of roots of $g$ having
valuation $r_h$. Factorization over $\overline{\Q}_p$ and multiplicativity
of the Gauss valuation give
\[
 \sum_{\substack{1\le h\le H\\r_h\le\rho}}\#R_h\,r_h
 \le v_\rho(g)\le C.
\]
Indeed, a root $\alpha$ of positive valuation contributes
$\min(v(\alpha),\rho)$ to $v_\rho(g)-v_0(g)$, and it belongs to at most
one of the $R_h$. The remaining term $v_0(g)$ is nonnegative.
Consequently some admissible $h$ has
$\#R_h\,r_h\le\delta/(h+1)$.

If $\alpha\in R_h$, then $1+\alpha$ is a unit. For this $\alpha$,
\[
 \sum_{\zeta\text{ primitive of order }p^h}
      \min\{v(\zeta-1-\alpha),1\}\le h+1.
\]
To see this, use the polynomial
$q(Y)=(Y+1+\alpha)^{p^h}-1$. Its coefficient of $Y$ has valuation $h$,
so $v_1(q)\le h+1$. The numbers $\zeta-1-\alpha$ in the displayed sum
are among its roots; the root-factorization formula for $v_1(q)$ gives
the inequality. Sum over $R_h$ and average over the
$\varphi(p^h)=1/r_h$ primitive roots. Some $\zeta$ then satisfies
\[
 \sum_{\alpha\in R_h}\min\{v(\zeta-1-\alpha),1\}
 \le\#R_h\,r_h(h+1)\le\delta<1.
\]
No summand can equal $1$, so the minima can be removed and none of the
differences vanishes. Roots at other valuations contribute exactly their
Gauss-valuation terms. Hence
\[
 v(g(\zeta-1))
 =v_{r_h}(g)+\sum_{\alpha\in R_h}
       \bigl(v(\zeta-1-\alpha)-r_h\bigr)
 \le v_{r_h}(g)+\delta.
\]
This is the source's root-counting and averaging argument.
\end{proof}

Fix an upper bound $s_0\ge1$ on the number of variables, a positive
rational $\rho$, a total valuation budget $B\ge0$, and an allowed loss
$\epsilon>0$. Define
\begin{equation}
\label{eq:torsion-parameters}
 h_* = \min\{h\ge1:r_h\le\rho\},\qquad
 \delta=\frac{\min(\epsilon,1)}{s_0+2}.
\end{equation}
Choose $H$ so that
\begin{equation}
\label{eq:harmonic-budget}
 B+1<\delta\sum_{h=h_*}^{H}\frac1{h+1}.
\end{equation}
A completely explicit sufficient choice is
\begin{equation}
\label{eq:explicit-torsion-height}
 H=\left\lceil(h_*+1)
             \exp\!\left(\frac{B+1}{\delta}\right)\right\rceil.
\end{equation}
The integral comparison
$\sum_{h=h_*}^{H}(h+1)^{-1}\ge
\log((H+2)/(h_*+1))$ proves the strict inequality in
\eqref{eq:harmonic-budget}. This is a sufficient bound extracted from
the proof, not an asserted optimal bound. The maximum order $p^H$ can
be roughly double exponential in $1/\epsilon$; the result is primarily
a structural existence theorem.

\begin{proposition}[Finite polynomial selection with positive-excess control]
\label{prop:polynomial-excess}
With these parameters and $1\le s\le s_0$, let $P_i$ be a finite
family in $\overline{\Z}_p[T_1,\ldots,T_s]$. Suppose that algebraic
witness points $x_i$ satisfy
\[
 v(x_{ij})\ge\rho,\qquad
 t_i=v(P_i(x_i))<\infty,\qquad\sum_i t_i\le B.
\]
There is one tuple $\zeta=(\zeta_1,\ldots,\zeta_s)$ such that
\[
 \zeta_j^{p^H}=1,\qquad r_H\le v(\zeta_j-1)\le\rho,
 \qquad
 \sum_i\bigl(v(P_i(\zeta-1))-t_i\bigr)_+\le\epsilon.
\]
In particular all the selected values are nonzero.
\end{proposition}

\begin{proof}
Replace the coordinates one at a time, using the same replacement for
all polynomials. At a given step, free the coordinate being replaced
in each polynomial and freeze its other coordinates at its own current
point. This gives integral one-variable polynomials $g_i$. If $c_i$
are the current value valuations, their sum is at most $B+1$: after
$k\le s_0$ steps the total increase will be at most $k\delta\le1$.
The coordinate about to be replaced still has valuation at least
$\rho$, so
\[
 v_\rho\left(\prod_i g_i\right)
 =\sum_i v_\rho(g_i)\le\sum_i c_i\le B+1.
\]
Apply Lemma~\ref{lem:one-coordinate-torsion} to the product and let
$a_i=v(g_i(\zeta-1))$, $b_i=v_{r_h}(g_i)$. Its conclusion is
\[
 a_i\ge b_i,\qquad c_i\ge b_i,\qquad
 \sum_i a_i\le\sum_i b_i+\delta.
\]
The inequality $c_i\ge b_i$ uses $r_h\le\rho$. All these quantities
are finite. Therefore
\[
 \sum_i(a_i-c_i)_+
 \le\sum_i(a_i-b_i)\le\delta.
\]
In particular the total valuation grows by at most $\delta$, justifying
the inductive bound used above. The positive part of a sum of changes
is at most the sum of their positive parts. After all $s$ coordinates,
the total positive excess is at most $s\delta\le\epsilon$.

Every selected coordinate is a primitive root of an order $p^h$ with
$h_*\le h\le H$. This gives both stated radius bounds and the common
order bound. The estimates in the source's
\texttt{exists\_ellVal\_eval\_update\_le\_add} prove precisely the
three inequalities above; its final wrapper keeps the individual
conclusion and discards the stronger aggregate information.
\end{proof}

\subsection{Completion and infinitely many simultaneous constraints}

\begin{theorem}[Uniform torsion selection for bounded power series]
\label{thm:analytic-torsion-selector}
Keep the parameters in \eqref{eq:torsion-parameters}--
\eqref{eq:harmonic-budget}. Let $I$ be any index set, and let
\[
 F_i\in\mathcal O_{\C_p}[[T_1,\ldots,T_s]],\qquad
 x_i\in\C_p^s,\qquad 1\le s\le s_0.
\]
Assume
\[
 v(x_{ij})\ge\rho,\qquad
 t_i=v(F_i(x_i))<\infty,\qquad
 \sup_{J\subset I,\,J\text{ finite}}\sum_{i\in J}t_i\le B.
\]
Then a single tuple of roots of unity satisfies
\[
 \zeta_j^{p^H}=1,\qquad r_H\le v(\zeta_j-1)\le\rho,
 \qquad
 \boxed{\sup_{J\subset I,\,J\text{ finite}}
 \sum_{i\in J}\bigl(v(F_i(\zeta-1))-t_i\bigr)_+\le\epsilon.}
\]
Consequently every $F_i(\zeta-1)$ is nonzero and
\[
 v(F_i(\zeta-1))\le t_i+\epsilon,\qquad
 \sup_{J\text{ finite}}\sum_{i\in J}v(F_i(\zeta-1))\le B+\epsilon.
\]
The bound $H$ is independent of the cardinality of $I$, the complexity
of the series, and the coefficient fields of the series and witnesses.
\end{theorem}

\begin{proof}
Let $G_H$ be the finite set of tuples satisfying the two order and radius
conditions. It is nonempty because $H\ge h_*$. A coordinate has
$p^H-p^{h_*-1}$ possible values, so
\[
 |G_H|=(p^H-p^{h_*-1})^s.
\]

First fix a finite $J\subset I$. Choose
$M>B+\epsilon+\rho+1$ and a nonnegative integer $N$ with
$(N+1)r_H>M$. Truncate each $F_i$, $i\in J$, at total degree $N$,
obtaining $P_i$. At each witness and every point of $G_H-1$, all
coordinates have valuation at least $r_H$. Thus every omitted
monomial has valuation at least $(N+1)r_H$. There are only finitely
many monomials of each total degree; their valuations tend uniformly
to infinity as the degree grows. Completeness of $\C_p$ and its
nonarchimedean inequality therefore show that the series converge
there and
\begin{equation}
\label{eq:uniform-torsion-tail}
 v(F_i(y)-P_i(y))\ge(N+1)r_H>M
\end{equation}
at every relevant point $y$.

There are now finitely many polynomial coefficients and witness
coordinates. By density of $\overline{\Q}_p$ in $\C_p$, approximate
every coefficient of the $P_i$ by an integral algebraic number, with
error valuation greater than $M$, and call the resulting polynomials
$Q_i$. Approximate each witness coordinate similarly by $x'_{ij}$.
As $M>\rho$, the new witnesses still satisfy $v(x'_{ij})\ge\rho$.

Integral polynomials are $1$-Lipschitz for the maximum norm on the unit
polydisc. Indeed, telescope the difference of two products of integral
coordinates and then sum the monomials. Thus replacing these witness
coordinates changes polynomial values by valuation greater than $M$.
Replacing coefficients has the same precision everywhere on the
unit polydisc. Together with \eqref{eq:uniform-torsion-tail} this gives
\[
 v(Q_i(x'_i)-F_i(x_i))>M,
 \qquad v(Q_i(x'_i))=t_i,
\]
because $t_i\le B<M$.

Apply Proposition~\ref{prop:polynomial-excess} to the $Q_i$, $i\in J$.
It gives $\zeta\in G_H$ with positive-excess budget at most $\epsilon$.
Each selected polynomial value has valuation at most
$t_i+\epsilon\le B+\epsilon<M$. Its difference from
$F_i(\zeta-1)$ has valuation greater than $M$, by coefficient
approximation and the uniform tail bound. The ultrametric inequality
then gives the exact equality
\[
 v(F_i(\zeta-1))=v(Q_i(\zeta-1)).
\]
The desired budget follows for this finite $J$, without consuming any
additional error allowance.

For an arbitrary $I$, let $S_J\subset G_H$ consist of the tuples
satisfying the excess bound for the finite set $J$. Every $S_J$ is
nonempty. For finitely many sets $J_1,\ldots,J_k$,
\[
 S_{J_1\cup\cdots\cup J_k}\subseteq
 S_{J_1}\cap\cdots\cap S_{J_k},
\]
since all excess summands are nonnegative. The family $(S_J)$ has the
finite intersection property inside a finite set; its total
intersection is therefore nonempty. A tuple in that intersection
proves the theorem. Each singleton constraint rules out an infinite
valuation, and summing $v(F_i(\zeta-1))\le
t_i+(v(F_i(\zeta-1))-t_i)_+$ gives the final aggregate bound.
\end{proof}

\begin{remark}[What the completion step adds]
Over one fixed finite extension $K/\Q_p$, with ramification index $e$,
a series in $\mathcal O_K[[T]]$ whose constant term is not a unit has
valuation at every witness of depth $\rho$ at least
$\min(\rho,1/e)$. The total budget then permits only finitely many
such series; a unit constant term is nonzero throughout the open unit
polydisc. Thus allowing $\mathcal O_{\C_p}$ coefficients, or algebraic
coefficient fields with unbounded ramification, makes the
infinite-family conclusion materially stronger. The finite-family
analytic result already applies over a fixed $K$.
\end{remark}

\subsection{Matrices, cokernels, and the precise arithmetic application}

\begin{corollary}[Simultaneous inverse and denominator control]
\label{cor:torsion-matrices}
Let $A_i(T)$ be square matrices, of possibly different finite sizes,
with entries in $\mathcal O_{\C_p}[[T_1,\ldots,T_s]]$. Suppose the
witness values
$t_i=v(\det A_i(x_i))$ obey the hypotheses of
Theorem~\ref{thm:analytic-torsion-selector}. Then a single tuple in
$G_H$ makes all the matrices invertible and satisfies
\[
 \|A_i(\zeta-1)^{-1}\|_{\max}\le p^{t_i+\epsilon}.
\]
The same positive-excess and total budgets hold for their determinant
valuations.
\end{corollary}

\begin{proof}
Apply the theorem to $F_i=\det A_i$. Adjugate entries are integral at
the selected point, and
$A_i^{-1}=\operatorname{adj}(A_i)/\det A_i$.
\end{proof}

If the coefficients belong to a fixed $\mathcal O_K$, the evaluations
belong to the finite extension $L=K(\mu_{p^H})$. Smith normal form
over the discrete valuation ring $\mathcal O_L$ gives the exact identity
\[
 v(\pi_L)\,\operatorname{length}_{\mathcal O_L}
      \operatorname{coker}A_i(\zeta-1)
 =v(\det A_i(\zeta-1)).
\]
Thus the theorem controls normalized lengths, not merely
nonvanishing. For a rectangular presentation with $r$ rows, choose an
$r\times r$ minor that is nonzero at the witness and apply the theorem
to that minor. Its cokernel surjects onto the cokernel of the full
presentation, so its determinant valuation bounds the latter's
normalized length above.

This supplies a concrete transfer problem: construct integral analytic
presentation matrices for a family of Selmer modules or
projection-denominator modules, and prove the nonzero-witness budget.
Once those arithmetic inputs are supplied, one finite-order character
specialization can control many modules at once. The theorem does not
construct such a presentation or establish the budget for an arbitrary
elliptic curve.

\subsection{Boundaries that prevent false conjecture proofs}

\paragraph{Arbitrary analytic functions are too general.}
The function
\[
 \log(1+T)=\sum_{n\ge1}(-1)^{n+1}\frac{T^n}{n}
\]
converges for $v(T)>0$. It has a nonzero value at $T=p^2$, of
valuation $2$, but vanishes at every $T=\zeta-1$ with $\zeta$ of
$p$-power order. Its coefficient valuations $-v_p(n)$ are unbounded
below. No fixed denominator makes this an integral bounded series.
The coefficient hypothesis is therefore essential.

\paragraph{Infinitely many separate bounds do not give a total budget.}
For $h\ge1$, take
\[
 F_h(T)=\Phi_{p^h}(1+T)\in\Z_p[T],\qquad x_h=p^2.
\]
The identity
$\Phi_{p^h}(u)=(u^{p^h}-1)/(u^{p^{h-1}}-1)$ and the elementary
lifting-the-exponent calculation for $u=1+p^2$ show that
$v(F_h(p^2))=1$ for every $h$, including $p=2$.
Nevertheless each nontrivial $p$-power root annihilates one $F_h$.
A separate bound $t_h\le1$ cannot replace the finite total budget.

\paragraph{The order may vary inside the finite grid.}
A single polynomial $\Phi_{p^h}(1+T)$ rules out every root of exact
order $p^h$, while having a finite witness valuation. Thus the result
does not allow an arbitrary prescribed exact order or conductor.
Likewise the product over $1\le h\le H$ vanishes at every nontrivial
$p$-power root of order at most $p^H$ and has witness valuation $H$.
Dependence on the valuation budget cannot be discarded. No density
statement about successful characters is proved.

\paragraph{Small normalized length does not automatically mean zero.}
The character and its coefficient field change as $\epsilon$ changes.
In a ramified field the smallest positive normalized length is
$v(\pi_L)=1/e_L$, and $e_L$ may grow rapidly with $H$. Consequently
producing, for every $\epsilon>0$, a different specialization of
normalized length at most $\epsilon$ does not by itself produce a
zero Selmer module. One must also control the coefficient field or
prove a suitable uniform discreteness statement. This prevents a
spurious deduction of Selmer vanishing from arbitrarily small losses.

\subsection{A structural obstruction for the unmodified elliptic Tate module}

The source's general twisting theorem also has a restriction that
cannot be ignored. The structure \texttt{New.Selmer.SelmerDatum} in
\texttt{SelmerData.lean} requires an inertia element $\gamma_*$ at a
place away from $p$ that acts by a scalar $\eta$, with
$\eta^2-1$ a unit. This hypothesis is incompatible with an unmodified
rank-two elliptic Tate module.

\begin{proposition}[Scalar-inertia obstruction]
\label{prop:scalar-inertia-obstruction}
Let $\rho$ be a rank-two representation over a nonzero integral
domain. If $\det\rho$ is trivial on inertia away from $p$, then no
element of that inertia can act as $\eta I_2$ with $\eta^2-1$ a unit.
\end{proposition}

\begin{proof}
Taking determinants gives $1=\det(\eta I_2)=\eta^2$.
Hence $\eta^2-1=0$, which is not a unit.
\end{proof}

For an elliptic curve, the Weil pairing identifies
$\bigwedge^2T_pE$ with $\Z_p(1)$. The determinant is therefore the
cyclotomic character, which is unramified at places of residue
characteristic different from $p$. The proposition applies, as it
also does to its Tate twists and dual. A preliminary finite-order
character twist can change the determinant; relating the resulting
local conditions and Selmer groups back to the desired arithmetic
problem is additional work.

The most promising connection is therefore to determinant and
specialization methods in $p$-adic families, including the missing
cases of refined Kurihara and Iwasawa statements. Existing work already
proves simultaneous nonvanishing for many finite families. The
specific added tool here is a uniform finite grid and a controlled
total valuation loss for bounded analytic data. It does not establish
an Iwasawa determinant identity, create a nonzero Euler-system class
from no initial witness, remove a functional-equation forced zero, or
prove a positive proportion of nonvanishing twists.

\subsection{Relation to existing nonvanishing and rigidity theorems}

Kriz--Nordentoft already develop horizontal \(p\)-adic \(L\)-functions,
refined maximum-modulus principles, and simultaneous nonvanishing under
specified arithmetic hypotheses~\cite{kn}. Serban's infinitesimal
multiplicative Manin--Mumford theorem addresses rigidity of torsion
intersections for formal schemes over finite extensions of
\(\Q_p\)~\cite{serban}. These results make the surrounding area an
established research program. The result proved here is the particular
finite-grid, total-positive-loss statement for bounded analytic data;
this analysis does not establish that it is absent from the entire
literature.

The matrix corollary suggests a disciplined route to a genuinely new
arithmetic theorem: identify a specific Selmer presentation or a
determinantal ideal in a \(p\)-adic family, normalize its coefficients
integrally, and produce finite nonzero witnesses with a total budget.
The selector then supplies a simultaneous specialization. The witness
step and the passage from a bounded normalized length to the exact
desired arithmetic assertion must be proved separately.
Section~\ref{sec:counterexample} gives another precise boundary on
overly strong specialization rules.


\clearpage
\section{A concrete nonabelian transfer: an \texorpdfstring{$A_5$}{A5} detector}
\label{sec:a5-transfer}

The earlier packets already propose $S_3$ and completely split-prime targets.
Here is a more demanding test case, together with a fully proved reduction to
one specific missing estimate.  The reduction uses the repository's one-sided
Landau mechanism and the now available RH for $\zeta$.  It does not assume
Artin's holomorphy conjecture.  The finite group calculations are independently
checked by the accompanying standard-library Python program
\texttt{a5\_detector.py}.  The arithmetic prime-sum estimate below is not proved.

Let $L/\Q$ be a Galois extension with group $A_5$.  Choose subgroups
$A_4,D_5,C_5$ of orders $12,10,5$, where $D_5$ denotes the dihedral group of
order $10$, and put
\[
 K_5=L^{A_4},\qquad K_6=L^{D_5},\qquad K_{12}=L^{C_5}.
\]
The subscripts are the degrees over $\Q$.  Write $\pi_j$ for the permutation
character of $A_5$ on the cosets defining $K_j$.

\begin{center}
\begin{tabular}{c|rrrrr}
Conjugacy class&$1$&$2A$&$3A$&$5A$&$5B$\\\hline
Class size&$1$&$15$&$20$&$12$&$12$\\
$\pi_5$&$5$&$1$&$2$&$0$&$0$\\
$\pi_6$&$6$&$2$&$0$&$1$&$1$\\
$\pi_{12}$&$12$&$0$&$0$&$2$&$2$\\\hline
$W=\pi_5+\pi_6+\pi_{12}-2$&$21$&$1$&$0$&$1$&$1$
\end{tabular}
\end{center}

\begin{lemma}[Two exact character identities]
\label{lem:a5-identities}
The function $W$ is nonnegative and has group average $1$.  The regular
character satisfies
\[
 \operatorname{reg}_{A_5}=4\pi_5+2\pi_6+3\pi_{12}-8.
\]
Consequently the following identity of meromorphic functions holds:
\begin{equation}
 \zeta_L(s)\zeta(s)^8
 =\zeta_{K_5}(s)^4\zeta_{K_6}(s)^2\zeta_{K_{12}}(s)^3.
 \label{eq:a5-zeta-identity}
\end{equation}
\end{lemma}

\begin{proof}
For a subgroup $H$, the coset character at $g$ is
$\#\{x\in A_5:x^{-1}gx\in H\}/|H|$.  This gives the table, and the
regular-character identity follows in each of its five columns.  The mean of
$W$ is $(21+15+12+12)/60=1$.  Artin formalism sends a permutation character on
$G/H$ to $\zeta_{L^H}$, sends sums to products, and sends the regular character
to $\zeta_L$.  Applying these identities first in $\Re s>1$, where Euler
products converge absolutely, and then continuing meromorphically proves
\eqref{eq:a5-zeta-identity}.  This use of Artin formalism does not assert the
holomorphy of individual nontrivial Artin factors; see, for example,
\cite{gun-hazra-sahu,ltz2024}.
\end{proof}

\subsection{A positive Euler coefficient at every prime power}

Define
\begin{equation}
 \mathcal F(s)=
 \frac{\zeta_{K_5}(s)\zeta_{K_6}(s)\zeta_{K_{12}}(s)}{\zeta(s)^2},
 \qquad
 D(s)=-\frac{\mathcal F'(s)}{\mathcal F(s)}.
 \label{eq:a5-detector-function}
\end{equation}
For a rational prime $p$, choose an inertia subgroup $I_p$ and a lift $g_p$
of its Frobenius element in the decomposition group.  Set
\begin{equation}
 a(p^k)=\log p\,\frac1{|I_p|}\sum_{h\in I_p}W(g_p^k h)
 \quad(k\geq1),
 \qquad a(n)=0\quad\hbox{otherwise}.
 \label{eq:a5-coefficients}
\end{equation}
The expression is independent of the choices.  Every summand in
\eqref{eq:a5-coefficients} is nonnegative, and
\[
 0\leq a(n)\leq21\Lambda(n),\qquad
 D(s)=\sum_{n\geq1}\frac{a(n)}{n^s}\quad(\Re s>1).
\]
Indeed, the projection onto inertia invariants is
$|I_p|^{-1}\sum_{h\in I_p}\rho(h)$, and its trace after applying
$\rho(g_p)^k$ is the average in \eqref{eq:a5-coefficients}.  The same identity
holds for a virtual character by linearity.  Expanding the local logarithmic
derivative gives the displayed Dirichlet series.  This is also the local
definition in \cite{gun-hazra-sahu}.  Thus neither ramified primes nor prime
powers are being discarded.  At an unramified prime the formula simplifies to
$a(p^k)=W(g_p^k)\log p$; using only $W(g_p)$ at all powers would be incorrect.

For example, at unramified primes the local factors of $\mathcal F$, with
$t=p^{-s}$, are
\[
\begin{array}{c|c}
1&(1-t)^{-21}\\
2A&(1-t)^{-1}(1-t^2)^{-10}\\
3A&(1-t^3)^{-7}\\
5A,5B&(1-t)^{-1}(1-t^5)^{-4}.
\end{array}
\]
These follow by computing the logarithmic derivative using the values of
$W(g^k)$, or directly from the three permutation actions.

\subsection{The precise missing estimate and its implication}

Each Dedekind zeta function has a simple pole at $1$, so $\mathcal F$ has a
simple pole there.  Let $b_{\mathcal F}$ be the constant term defined by
\[
 D(s)=\frac1{s-1}+b_{\mathcal F}+O(s-1),
 \qquad c_{\mathcal F}=b_{\mathcal F}-1,
\]
and define, for $x\geq1$,
\begin{equation}
 R_W(x)=\sum_{n\leq x}a(n)\left(\frac1n-\frac1x\right)
           -\log x-c_{\mathcal F}.
 \label{eq:a5-remainder}
\end{equation}
This specifies the constant without an unknown normalization.  In terms of
Euler--Kronecker constants it is
$c_{\mathcal F}=-1-\gamma_{K_5}-\gamma_{K_6}-\gamma_{K_{12}}+2\gamma$.

\begin{theorem}[A single nonabelian arithmetic target]
\label{thm:a5-criterion}
Assume RH for the Riemann zeta function.  Suppose that for every
$\varepsilon>0$ there are $C_\varepsilon,X_\varepsilon>0$ such that
\begin{equation}
 R_W(x)\leq C_\varepsilon x^{-1/2+\varepsilon}
 \qquad(x\geq X_\varepsilon).
 \tag{$A_5$-bound}\label{eq:a5-missing-bound}
\end{equation}
Then $\zeta_L$, $\zeta_{K_5}$, $\zeta_{K_6}$, and $\zeta_{K_{12}}$ satisfy
their Riemann hypotheses.  No holomorphy hypothesis for the individual
irreducible Artin $L$-functions is needed.
\end{theorem}

\begin{proof}
Fix $1/2<\beta<1$.  Choose $0<\varepsilon<\beta-1/2$; by increasing the
threshold, \eqref{eq:a5-missing-bound} implies
$R_W(x)\leq x^{\beta-1}$.  Put $E(v)=R_W(e^v)$ for $v\geq0$, and extend it
by zero for $v<0$.  The coefficient bound above
implies $E(v)=O(1+v^2)$ for $v\geq0$, so in particular $E$ is measurable and
has the growth allowed by \texttt{New.IsPolar}.  Direct integration in
$\Re s>1$ gives
\begin{align}
 \int_0^\infty E(v)e^{(1-s)v}\,dv
 &=\int_1^\infty R_W(x)x^{-s}\,dx\notag\\
 &=\frac{D(s)}{s(s-1)}-\frac1{(s-1)^2}
                         -\frac{c_{\mathcal F}}{s-1}.
 \label{eq:a5-transform}
\end{align}
To check the first term, integrate each ramp from $n$ to $\infty$:
\[
 \int_n^\infty\left(\frac1n-\frac1x\right)x^{-s}\,dx
 =\frac{n^{-s}}{s(s-1)}.
\]
At $s=1$ the Laurent expansion of the first term in
\eqref{eq:a5-transform} is
$(s-1)^{-2}+(b_{\mathcal F}-1)(s-1)^{-1}+O(1)$, so the two prescribed
subtractions cancel both poles.  Elsewhere in $\Re s>0$ the transform is
meromorphic with at most simple poles.  If $m(\rho)$ is the signed order of
$\mathcal F$ at $\rho\ne1$, its residue is
\begin{equation}
 r(\rho)=\frac{m(\rho)}{\rho(1-\rho)}.
 \label{eq:a5-residue}
\end{equation}
The sign here is essential: $-\mathcal F'/\mathcal F$ has residue $-m(\rho)$.

RH ensures that the denominator $\zeta(s)^2$ in
\eqref{eq:a5-detector-function} has no zero in $\Re s>1/2$.  Away from $1$,
all three numerator factors are holomorphic, whence $m(\rho)\geq0$ in that
half-plane.  For real $\beta<u<1$, \eqref{eq:a5-residue} is therefore
nonnegative.  At $u=1$ the corrected transform is holomorphic, and for real
$u>1$ the Euler products give $m(u)=0$.  The one-sided bound and these real
residue signs are exactly the input of
\texttt{New.IsPolar.landau} in \texttt{Lib/Analytic/Landau.lean}.
It follows that $r(\rho)=0$ for $\Re\rho>\beta$.  Thus $\mathcal F$ has no
zero there.  Since $\beta>1/2$ was arbitrary, this excludes all its zeros
with real part greater than $1/2$.

In this half-plane, away from $1$, the order of $\mathcal F$ is the sum of
the three nonnegative orders of its numerator zeta functions.  Their zeros
are therefore all excluded individually.  Equation
\eqref{eq:a5-zeta-identity} then excludes zeros of $\zeta_L$ as well.
The functional equations of the four Dedekind zeta functions reflect any
nontrivial zero with real part below $1/2$ to one above $1/2$.  This proves
the conclusion.
\end{proof}

The theorem is a sufficient criterion, not a proof of
\eqref{eq:a5-missing-bound}.  Its useful feature is the separation of tasks:
the analytic implication, treatment of real zeros, ramification, prime powers,
and denominator poles are complete.  A new short/enough/long-forms construction
must now reproduce the coefficients $W(\operatorname{Frob}_p^k)$ and prove
the displayed upper bound.  The source's current congruence and quadratic
weights do not supply it.

\subsection{A scoped optimality statement}

Consider only the normalized three-field family
\[
 W_{a,b,c}=a\pi_5+b\pi_6+c\pi_{12}+(1-a-b-c)\cdot1,
 \qquad a,b,c\geq1.
\]
The normalization gives mean $1$; the lower bounds require each of the three
Dedekind factors to be detected with coefficient at least $1$.  Its identity
value is
\[
 W_{a,b,c}(1)=1+4a+5b+11c\geq21.
\]
Thus any nonnegative member has peak at least $21$.  The table proves that
$a=b=c=1$ attains peak $21$.  This is an exact optimum in this declared family,
also for real coefficients.  It is not an assertion of global optimality
among all subfields, all kernels, or all analytic reductions.  The direct
regular-character detector has peak $60$ at the same mean $1$.

\subsection{Why fixed abelian congruence weights cannot replace this input}

\begin{theorem}[An exact mean-square obstruction]
\label{thm:a5-blindness}
Let $B/\Q$ be any finite abelian extension with group $H$, and let
$f:H\to\C$ be arbitrary.  For primes unramified in $LB$, write $g_p,h_p$
for the two Frobenius components.  Then
\[
 \lim_{X\to\infty}\frac1{\pi(X)}
 \sum_{\substack{p\leq X\\p\text{ unramified in }LB}}
 |W(g_p)-f(h_p)|^2
 =7+\frac1{|H|}\sum_{h\in H}|f(h)-1|^2\geq7.
\]
Equality holds for $f\equiv1$.  In particular, no fixed finite collection of
Dirichlet characters can represent the weight $W$ at all unramified primes.
\end{theorem}

\begin{proof}
$A_5$ has no nontrivial abelian quotient.  Therefore $L\cap B=\Q$ and
$\operatorname{Gal}(LB/\Q)=A_5\times H$.  Chebotarev turns the prime average
into the uniform group average.  The table gives
\[
 \frac1{60}\sum_{g\in A_5}W(g)=1,
 \qquad
 \frac1{60}\sum_{g\in A_5}(W(g)-1)^2
 =\frac{400+20}{60}=7.
\]
Expanding $|W-f|^2=|(W-1)+(1-f)|^2$, the mixed term vanishes after the
$A_5$ average; the displayed identity follows.  Any finite collection of
Dirichlet characters factors through one finite cyclotomic, hence abelian,
extension, so the final assertion follows as well.
\end{proof}

This obstruction concerns coefficient representations through a fixed finite
abelian extension.  It is not a logical independence theorem, does not rule
out a new argument using RH, and does not cover moduli changing with the
prime-sum scale.  Such changing moduli would in turn require new uniform
estimates, whereas the source's final bounds fix its modulus first.  The
general phenomenon that $A_5$ splitting behavior is not determined by
congruences is classical; the exact variance $7$ above is the calculation for
this particular detector.

\subsection{A completely specified field}

One can take $L$ to be the splitting field of
\[
 P(X)=X^5+20X+16.
\]
Here is a short Galois certificate, reproduced exactly by the companion
program.  The polynomial discriminant is
$1\,024\,000\,000=32\,000^2$, so the Galois group lies in $A_5$.
Modulo $3$, $P=X^5-X+1$ has no root.  Its remainders modulo the three monic
irreducible quadratics $X^2+1$, $X^2+X+2$, and $X^2+2X+2$ are respectively
$1$, $X+1$, and $X+1$.  A reducible quintic has a factor of degree at most
$2$, so $P$ is irreducible modulo $3$.  Modulo $7$,
\[
 P=(X+2)(X+3)(X^3+2X^2-2X-2),
\]
and the cubic has no root.  Thus the Galois group contains elements of orders
$5$ and $3$.  Its order divides $60$ and is divisible by $15$, leaving
$15,30,60$.  A group of order $15$ is cyclic by Sylow's theorems and cannot
embed in $S_5$; an index-two subgroup of $A_5$ would give a nontrivial map to
$C_2$, impossible since $A_5$ is generated by elements of order $3$.
Hence the group has order $60$ and is $A_5$.

This identifies a fixed, genuinely nonabelian zeta-function target.  No
claim is made that its RH has been proved by the present calculation.  Nor
does this report claim priority for the polynomial or Artin factorization.
The contribution is the explicit positive detector, its exact analytic
certificate and scoped optimization, and the quantified obstruction to
replacing its nonabelian coefficients by fixed congruence data.

\clearpage
\section{An extension of the exponent-six exclusions}
\label{sec:genus-extension}

The earlier packet \texttt{grh\_transfers.zip} ends by asking whether
smaller subfields of a genus field, together with ramified ideals of
small norm, can improve its exponent-six exclusions. The following
argument answers that specific question. It uses the exported Dirichlet
GRH and an explicit prime bound; it does not extend the theta-lift
construction to a new family of $L$-functions. The mathematical deduction
is new relative to the five packets reviewed here. No publication
priority or Lean verification is claimed.

\subsection{A general certificate}

Let $D<-4$ be a negative fundamental discriminant, put $x=|D|$, and write
$K=\mathbb Q(\sqrt D)$ and $A=\operatorname{Cl}(K)$. A ramified ideal
will mean a squarefree product of the unique prime ideals above rational
primes dividing $D$. Let $t=\omega(x)$, the number of those rational primes.
The genus field $G$ satisfies
\[
 [G:\mathbb Q]=2^t,\qquad
 \operatorname{Gal}(G/K)\simeq A/A^2.
\]
It is multiquadratic over $\mathbb Q$ and unramified over $K$ at every
finite prime. These are the standard genus-field facts used in the
previous packet.

\begin{theorem}[Intermediate genus-field certificate]
\label{thm:genus-intermediate}
Suppose the exponent of $A$ divides $2m$, where $m$ is a positive odd
integer. Let $H\subseteq A[2]$ have dimension $k$ over $\mathbb F_2$.
Assume that each class in $H$ has a ramified ideal representative of norm
at most $R$. There is an abelian field $F/\mathbb Q$ containing $K$, of
degree $n=2^{t-k}$ and absolute discriminant $x^{n/2}$, such that every
unramified rational prime $p$ splitting completely in $F$ satisfies
\begin{equation}\label{eq:genus-norm}
                         x\le 4R p^m.
\end{equation}
Consequently the repository's Dirichlet GRH implies
\begin{equation}\label{eq:genus-universal}
 x\le 4R\bigl(2n\log x+\tfrac52n+5\bigr)^{2m}.
\end{equation}
Any sharper valid upper bound for the least completely split prime in
$F$ may replace the bound used in \eqref{eq:genus-universal}.
\end{theorem}

\begin{proof}
The map $A[2]\to A/A^2$ is an isomorphism. Indeed, $A^2$ has exponent
dividing the odd integer $m$, so it intersects $A[2]$ trivially. For
surjectivity, $a^m\in A[2]$ and $a^{1-m}$ is a square for every $a\in A$.
Regard $H$ as a subgroup of the genus Galois group, and put $F=G^H$.
Then $K\subseteq F$, $[F:\mathbb Q]=n$, and $F/K$ is unramified.
The discriminant tower formula gives
\[
                   |\Delta_F|=|D|^{[F:K]}=x^{n/2}.
\]
Moreover $F/\mathbb Q$ is abelian because $G/\mathbb Q$ is.

Let $p$ be as stated. Since $F$ contains $K$, the prime $p$ splits in
$K$ and is prime to $D$. Choose one prime $\mathfrak p$ of $K$ above $p$.
Its Artin symbol in $\operatorname{Gal}(G/K)$ lies in $H$, so
$[\mathfrak p]=h a^2$ for some $h\in H$ and $a\in A$. Raising to the
$m$th power gives $[\mathfrak p]^m=h$. Choose a ramified ideal $J$ in
class $h$ with $N(J)\le R$. Because $h^2=1$, the integral ideal
$\mathfrak p^mJ$ is principal, say $(\alpha)$.

The valuations of $(\alpha)$ at $\mathfrak p$ and
$\overline{\mathfrak p}$ are $m$ and zero: $J$ is supported only at
primes above $D$. Hence $\alpha\notin\mathbb Q$. Write
$\alpha=(a+b\sqrt D)/2$ with integers $a,b$ and $b\ne0$. Norms give
\[
 p^mN(J)=N_{K/\mathbb Q}(\alpha)
        =\frac{a^2+xb^2}{4}\ge\frac x4,
\]
proving \eqref{eq:genus-norm}. Notice that the norm cost is $R$, not
$R^m$: the ramified ideal is inserted \emph{after} taking the $m$th power.

The Dedekind zeta function of the abelian field $F$ factors into
Dirichlet $L$-functions, so the exported theorem supplies its RH.
Theorem 5.1 of Bach--Sorenson gives an unramified prime with identity
Frobenius and
\[
 p\le\bigl(4\log|\Delta_F|+\tfrac52n+5\bigr)^2.
\]
Substitution proves \eqref{eq:genus-universal}.
\end{proof}

The explicit prime theorem is in E. Bach and J. Sorenson,
\href{https://doi.org/10.1090/S0025-5718-96-00763-6}{\emph{Explicit bounds
for primes in residue classes}}, Math. Comp. 65 (1996), Theorem 5.1
and Table 3. Here it is used only for an abelian extension of
$\mathbb Q$, so no additional nonabelian GRH assumption is involved.

One may optimize the theorem by replacing $R$ with
\[
 R(H)=\max_{h\in H}\min\{N(J):J\text{ ramified and }[J]=h\},
\]
then varying $H$. The following explicit choice already gives a useful
improvement and avoids any search over subgroups.

\subsection{The rank issue, including the prime two}

Let $q_1<\cdots<q_t$ be the rational primes dividing $D$, with
ramified prime ideals $\mathfrak q_i$, and write
$P_j=q_1\cdots q_j$, with $P_0=1$.

\begin{lemma}\label{lem:ramified-independent}
Under the exponent hypothesis of Theorem~\ref{thm:genus-intermediate},
any selection of at most $t-2$ ramified prime classes is independent in
$A[2]$ and in $A/A^2$. Thus for $0\le k\le t-2$ the subgroup generated
by $[\mathfrak q_1],\ldots,[\mathfrak q_k]$ has dimension $k$ and has
ramified representatives of norm at most $P_k$.
\end{lemma}

\begin{proof}
Each ramified class has order dividing two. Suppose a nonempty
squarefree product $J$ of ramified ideals is principal, say $(\alpha)$.
Conjugation fixes $J$, so $\overline\alpha/\alpha$ is a unit of $K$.
Since $D<-4$, the units are $\{1,-1\}$. If
$\overline\alpha=\alpha$, then $\alpha$ is rational and its valuations
at all ramified primes are even, contrary to the definition of $J$.
Otherwise $\alpha$ is a rational multiple of $\sqrt D$.

Multiplication by a rational number changes every ramified valuation by
an even integer. At each odd prime dividing $D$, the valuation of
$\sqrt D$ is odd. At the prime above two its parity is the parity of
$v_2(D)$. A fundamental even discriminant has $v_2(D)=2$ or $3$.
Therefore a nonempty relation has support all $t$ ramified primes when
$D$ is odd or $v_2(D)=3$, and all $t-1$ odd ramified primes when
$v_2(D)=2$. No selection of at most $t-2$ can support a relation.
Finally, a class of order dividing two that lies in $A^2$ is trivial,
since $A^2$ has odd exponent. This proves independence also in the genus
quotient. Every class of the stated subgroup has a representative formed
from a subset of the first $k$ prime ideals, of norm at most $P_k$.
\end{proof}

\subsection{Four certificates in the unresolved exponent-six window}

\begin{theorem}[Additional necessary conditions in the census window]
\label{thm:genus-four}
Assume the repository's Dirichlet GRH. Let $D<-4$ be fundamental, let
$\operatorname{Cl}(\mathbb Q(\sqrt D))$ have exponent six, and suppose
\[
 3.1\cdot10^{20}<x=|D|\le2.5\cdot10^{25},\qquad t=\omega(x)\ge7.
\]
Then, with the products $P_j$ defined above,
\begin{equation}\label{eq:genus-four}
 P_{t-2}>1{,}200{,}000,\qquad
 P_{t-3}>75{,}000,\qquad
 P_{t-4}>850,\qquad
 P_{t-5}>14.
\end{equation}
\end{theorem}

\begin{proof}
Put $y=\log x$, so $47<y<59$. Apply
Theorem~\ref{thm:genus-intermediate} with $m=3$ and
$k=t-\log_2 n$, for $n=4,8,16,32$. The preceding lemma verifies the
rank in all four cases. Table 3 of Bach--Sorenson supplies
$\sqrt p\le M_n(y)$ with the following affine upper bounds:
\begin{center}
\begin{tabular}{rrrrr}
\toprule
$n$ & $k$ & $M_n(y)$ & Test value $R$ & Ratio at $y=47.2$\\
\midrule
4 & $t-2$ & $4.070y+6.38$ & $1{,}200{,}000$ & $<0.946745$\\
8 & $t-3$ & $6.516y+6.98$ & $75{,}000$ & $<0.937077$\\
16 & $t-4$ & $13.960y+6.44$ & $850$ & $<0.951538$\\
32 & $t-5$ & $27.920y+7.08$ & $14$ & $<0.977087$\\
\bottomrule
\end{tabular}
\end{center}
The last column bounds $4R M_n(47.2)^6/(3.1\cdot10^{20})$ and is
displayed only for readability; the decisive comparisons use rational
arithmetic.

For $n=8,16,32$, the logarithm of the field discriminant is
$(n/2)y$, lying respectively in $(188,236)$, $(376,472)$, and
$(752,944)$. The applicable Table-3 triples $(A,B,C)$ in
$\sqrt p\le A\log|\Delta_F|+Bn+C$ are
$(1.629,0.11,6.1)$ for $n=8$, and $(1.745,0.04,5.8)$ for $n=16,32$.
For $n=4$ and $2y\le100$, the triple is $(2.035,0.27,5.3)$.
For $2y\ge100$, the triple is $(1.527,0.17,6.4)$; its bound
$3.054y+7.08$ is below the displayed $4.070y+6.38$ when $y\ge50$.
Thus the stated envelopes cover the entire window.

For every row, $4R(ay+b)^6e^{-y}$ decreases when $y>6-b/a$, in
particular when $y>47$. Write $X_0=3.1\cdot10^{20}$. Since
$\log X_0<236/5$, the exact checks
\[
                        4R M_n(236/5)^6<X_0
\]
give $4R M_n(y)^6<x$ throughout $x>X_0$. If the corresponding $P_k$
were at most $R$, equation~\eqref{eq:genus-norm} would imply the reverse
inequality. This proves all four necessary conditions.
\end{proof}

The published EKN census and global bound are independent inputs to its
classification interpretation: Elsenhans--Kl\"uners--Nicolae,
\href{https://arxiv.org/pdf/1803.02056}{\emph{Imaginary quadratic number
fields with class groups of small exponent}}, Theorems 1 and 3,
Tables 1 and 2. They give 432 exponent-six fields through $X_0$ and,
under the quadratic RH supplied here, no exponent-six field above
$2.5\cdot10^{25}$. The preceding packet excludes the remaining cases
with $t\le6$ and those with $x/P^+(x)\le10^8$.
Theorem~\ref{thm:genus-four} therefore excludes additional families
from the remaining window. It does not enumerate all survivors or prove
completeness of the 432-field list without a family restriction.

\paragraph{A strict improvement over those stated conditions.}
The fundamental discriminant
\[
\begin{split}
D&=-8\cdot3\cdot5\cdot7\cdot11\cdot1000000007\cdot1000000009\\
 &=-9240000147840000582120
\end{split}
\]
has seven distinct ramified primes. It satisfies
$X_0<|D|<2.1\cdot10^{22}$ and
$|D|/P^+(|D|)=9240000064680>10^8$, so it passes the earlier packet's
stated numerical necessary conditions. But $P_5=2310<1{,}200{,}000$;
the new theorem rules out exponent six. This demonstrates a strict
enlargement of those exclusion families, not a claim that this individual
discriminant was previously difficult or unknown.

\paragraph{Verification.}
The accompanying \texttt{genus\_certificate.py} uses exact rational
arithmetic to verify all four endpoint comparisons and exponential-series
bounds for $47<\log X_0<47.2$ and $\log(2.5\cdot10^{25})<59$.
It certifies the two large example primes by trial division, and checks
fundamentality, the products and the prior numerical conditions.
The JSON output records the exact ratios. It does not replace the
mathematical proof or independently recheck the published census.

\section{A sharp boundary for maximal-order torsion selection}
\label{sec:counterexample}

The torsion theorem allows the successful conductor to vary within a
finite range. It is tempting to replace that flexibility by an
exact-order conclusion with no valuation loss. The following explicit
counterexample prevents one natural such strengthening over arbitrarily
ramified coefficient rings. It is a correction to that proposed principle,
not a claim to refute a theorem in the cited literature.

Let \(G=C_{p^m}\), \(m\ge2\), with generator \(g\), and let \(\zeta\)
be a primitive \(p^m\)-th root of unity. In the group ring
\(\mathcal O_{\Q_p(\zeta)}[G]\), define
\[
 \nu=\prod_{a=1}^{p-1}\bigl([g]-\zeta^a[0]\bigr).
\]
We use additive notation for \(G\), with identity \(0\).
A character evaluates \([g]\) at its value on \(g\); the trivial
character, denoted \(1\), is the augmentation. Put
\(r_m=1/\varphi(p^m)\).

\begin{proposition}
\label{prop:maximal-order-counterexample}
The augmentation and maximal-order evaluations satisfy
\[
 v\bigl(\nu(1)\bigr)=(p-1)r_m=p^{1-m}=:t<1,
 \qquad
 v\bigl(\nu(\chi)\bigr)\ge2t
\]
for every character \(\chi\) of exact order \(p^m\), where the second
valuation may be infinite. In particular no such character satisfies
\(v(\nu(\chi))\le v(\nu(1))\).
\end{proposition}

\begin{proof}
Each \(1-\zeta^a\), \(1\le a\le p-1\), has valuation \(r_m\), giving
the augmentation formula. Write \(\chi(g)=\zeta^b\) with \(p\nmid b\).
Exactly one \(a\in\{1,\ldots,p-1\}\) satisfies \(a\equiv b\pmod p\).
For every other \(a\), \(v(\zeta^b-\zeta^a)=r_m\). For the matching
\(a\), the difference is zero or has valuation at least \(p r_m\),
because \(p\mid b-a\). Summing gives
\[
 v\bigl(\nu(\chi)\bigr)\ge (p-2)r_m+p r_m=2(p-1)r_m=2t.
\]
\end{proof}

For example, \(p=2,m=2\) gives
\(\nu=[g]-i[0]\). Its augmentation has valuation \(1/2\);
the two primitive evaluations have valuations \(+\infty\) and \(1\).
Meanwhile \(v_p(|G[p]|)=1\), so the weaker condition
\(v(\nu(1))<v_p(|G[p]|)\) would not justify the false maximal-order,
no-loss conclusion.

The threshold \(1/p^{m-1}\) appearing in the refined maximum-modulus
principle of Kriz--Nordentoft, Proposition 2.11, is exactly where this
example sits~\cite{kn}. This observation explains why dropping
ramification-sensitive thresholds is dangerous. It does not address all
their hypotheses or settle their separate questions about groups of
unbounded exponent.

There is an important quantifier limit to the example: \(t=p^{1-m}\)
tends to zero as \(m\) grows. It rules out no loss, and also any specified
additive loss \(\epsilon<t\) for this example, but by itself does not rule
out a theorem that first fixes \(\epsilon>0\) and then chooses \(m\)
large enough. The separate polynomial
\(\Phi_{p^H}(1+T)\), discussed in Section~\ref{sec:torsion-transfer},
vanishes at every root of prescribed exact order \(p^H\) while having
witness valuation \(1\). That is the appropriate obstruction to a
universal exact-order selector for all bounded-budget polynomial inputs.

The script \file{code/torsion_checks.py} verifies these valuations
exactly for \(p=2,3,5,7\) and \(m=2,3,4\). It also uses cyclotomic
resultants to illustrate why varying the chosen order permits
simultaneous avoidance for a finite family.


\section{Research priorities and conjectures not yet resolved}
\label{sec:priorities}

\subsection{First priority: extract and formalize the analytic selector}

The strongest completed transfer is Theorem~\ref{thm:analytic-torsion-selector}.
It relies on a local root-counting argument and has a tractable
formalization path:
\begin{enumerate}
\item Strengthen the one-coordinate update lemma to retain
\(\sum_i(a_i-c_i)_+\), then telescope across coordinates.
\item Formalize uniform total-degree tails on a fixed torsion grid and
algebraic approximation of finitely many coefficients and witnesses.
\item Apply the finite-intersection argument on the finite grid.
\item Package determinant, maximal-minor, and normalized-length corollaries
as a reusable specialization API.
\end{enumerate}
No new global RH machinery is needed. The current 103-module import
closure is an engineering starting point; reducing incidental imports
would improve portability.

The result's explicit order bound is poor for computation. For
\(p=2,s_0=1,\rho=1/2,B=1,\epsilon=1/4\), the displayed sufficient
choice is
\[
 H=\lceil3e^{24}\rceil,\qquad \text{maximum torsion order }2^H.
\]
This is roughly \(H=7.95\cdot10^{10}\), before exponentiating again.
It proves uniform existence, not a feasible enumeration algorithm.
Improving the harmonic-pigeonhole bound for particular coefficient
classes or determinant structures is a concrete secondary problem.

\subsection{Second priority: attack the \texorpdfstring{\(A_5\)}{A5} arithmetic estimate}

The exact proposed nonabelian target is
\[
 R_W(x)\le C_\epsilon x^{-1/2+\epsilon},
 \qquad W=(21,1,0,1,1).
\]
The positive detector includes all prime powers and handles inertia.
The analytic conclusion, the pole normalization, the real residue signs,
the chosen field, and the limited character optimization have been proved.
They should be treated as finished infrastructure.

The unresolved task is to construct arithmetic objects whose index or
determinant encodes this \(W\), with matching main terms and errors
small enough to preserve the square-root saving. A reasonable first
attempt would make the degree-5, degree-6, and degree-12 permutation
realizations interact with a common lattice comparison. It would need
to prove integrality, sufficient rank, a compatible long-form estimate,
and local twist control. This is a proposed research route, not an
existing construction.

The mean-square lower bound \(7\) in
Theorem~\ref{thm:a5-blindness} gives a useful early failure test:
any proposal that uses only finitely many fixed congruence classes to
represent \(W\) exactly cannot work. A growing-modulus proposal must
supply new uniform estimates and explain how it retains the source's
scale. One cannot infer failure of every RH-based argument from this
limited obstruction.

The author announcement itself identifies extended RH as a possible
further direction~\cite{alpoge-remark}. It provides motivation, not
the missing nonabelian construction. Approximate Artin holomorphy
and zero-free results for families are also substantial existing work,
but they do not provide the fixed-field, full-critical-line statement
required by this target~\cite{ltz2024}.

\subsection{Third priority: finish a finite class-group classification}

Theorem~\ref{thm:genus-four} strengthens a specific gap in the prior
analysis. It is the closest direction to a complete further
classification theorem because a published finite global range and an
existing census are already available. A next computation should apply
the actual inequality
\[
 x\le4R(H)\,M_{2^{t-\dim H}}(\log x)^6
\]
for useful subgroups \(H\), rather than only the four rounded thresholds.
Optimizing ramified representatives can reduce \(R(H)\).
The remaining discriminants would then need a complete, independently
checkable enumeration or a further theoretical exclusion.

The present report does not count that remaining set, estimate its
computational cost, or verify all 432 published entries. Any claim that
the classification has now been completed would go beyond both the
proofs and the computations supplied here.

\subsection{Why BSD, Goldfeld, Vandiver, and general Iwasawa claims do not follow}

For BSD or Goldfeld-type nonvanishing, a selector must be connected to
the correct analytic or algebraic object. A nonzero initial witness,
integral interpolation or a presentation, compatibility of local
conditions, and any required density statement remain independent
inputs. The scalar-inertia obstruction shows that even the source's
packaged Selmer hypothesis fails for the unmodified elliptic Tate module.
A selector cannot remove a forced sign zero or create a nonzero
Euler-system class from an identically zero input.

For Vandiver-type or higher \(K\)-group vanishing statements, controlling
a denominator at varying ramified twists does not show that the desired
fixed component vanishes. Normalized lengths can decrease simply
because the coefficient-field ramification increases. One would need
a fixed-component comparison and a uniform positive lower gap for
nonzero lengths, or a different argument proving exact vanishing.
The previous reports already identify related barriers~\cite{prior}.

Thus these are candidate domains for the determinant machinery, not
conjectures resolved here. A useful follow-up should name one exact
representation, module, and specialization problem and verify its
hypotheses before claiming a new arithmetic consequence.

\section{Verification record and use of the archive}
\label{sec:repro}

\subsection{Claim and evidence ledger}

\begin{center}
\small
\begin{tabularx}{\textwidth}{@{}p{49mm}YY@{}}
\toprule
Claim & Evidence supplied & Limit\\
\midrule
Scope of upstream theorem &
Pinned \file{Main.lean}, \file{Statements.lean}, source hashes,
reported independent checking pages. &
No new local Lean build or kernel replay.\\
\addlinespace
Analytic torsion theorem &
Complete proof including positive excess, approximation, uniform tails,
and the finite-intersection step. &
Finite Python examples do not certify the infinite theorem.\\
\addlinespace
\(A_5\) transfer criterion &
Complete transform/residue proof and exact finite-group and polynomial
certificate. &
The prime-sum inequality remains a hypothesis.\\
\addlinespace
Genus-field exclusion theorem &
Complete class-group/norm proof, original prime-bound constants, exact
rational endpoint checks. &
Published census is an external input; no full survivor enumeration.\\
\addlinespace
Novelty relative to packets &
Inspection of all five prior archives, pinned manifest, detailed
comparison notes. &
No exhaustive literature-priority determination.\\
\bottomrule
\end{tabularx}
\end{center}

The exact finite checks have distinct purposes. The \(A_5\) program
constructs the group and coset actions, verifies the character identities,
local cycle factors, variance and the Galois certificate. The genus
program checks rational inequalities, table ranges, the demonstration
primes and discriminant. The torsion program checks illustrative finite
valuation computations and a harmonic bound. All decisive arithmetic in
the first two programs is exact. Approximate decimals in the JSON are
display aids, not the basis of comparisons.

\subsection{Contents and commands}

The archive contains \file{report.tex}, \file{report.pdf}, the included
section files, Python certificates, their JSON/CSV results,
source-audit and prior-report manifests, review notes, and a small
selection of upstream source files with their license notices.
The upstream excerpt is evidence, not a standalone buildable copy of
\repo.

Run \texttt{bash code/reproduce.sh} from the archive root to regenerate
the finite certificates and rebuild the PDF when the dependencies are
available. Python 3 is sufficient for the \(A_5\) and genus scripts;
the torsion illustrations also require SymPy. A standard LaTeX installation
with \file{latexmk} and \file{pdflatex} builds the report. The source audit
accepts the path to a separately cloned repository at the pinned revision.
The included SHA-256 manifest checks archive members; it does not attest
to their mathematical correctness.


\begingroup
\small
\begin{thebibliography}{99}
\addcontentsline{toc}{section}{References}

\bibitem{repo}
\texttt{lsernau/grh-lean}, readable and compact Lean formalizations,
revision \file{72fead102f663d5c7a407964a13dc0cf9d977207},
25 November 2026.
\url{https://github.com/lsernau/grh-lean/tree/72fead102f663d5c7a407964a13dc0cf9d977207}.
Specific paths and SHA-256 hashes are recorded in the accompanying audit.

\bibitem{prior}
\texttt{VladimirReshetnikov/ProveIt}, \file{docs/incoming},
revision \file{a92443f00663b1b8b3c233d17df0ad49dc4f2b2c}.
\url{https://github.com/VladimirReshetnikov/ProveIt/tree/a92443f00663b1b8b3c233d17df0ad49dc4f2b2c/docs/incoming}.
The five archive names, hashes, and internal report paths are in
\file{evidence/prior_reports_manifest.json}.

\bibitem{sutherland-check}
A. V. Sutherland, independent GRH formalization checking records.
\url{https://math.mit.edu/~drew/shortgrh/}.
Consulted as a report of external checks, not a local replay.

\bibitem{sonoda-check}
S. Sonoda, independent checking and endpoint-comparison report.
\url{https://shosonoda.github.io/lean-rh-alpoge/grh.html}.

\bibitem{rubinstein-check}
M. Rubinstein, compact GRH formalization checking report.
\url{https://mrubinst.github.io/grh-lean-compact-check/}.

\bibitem{alpoge-remark}
L. Alpoge, announcement and remarks on the GRH formalization, 2026.
\url{https://alpo.ge/remark.pdf}.

\bibitem{kn}
D. Kriz and A. C. Nordentoft, \emph{Horizontal \(p\)-adic \(L\)-functions},
arXiv:2310.20678, version 3, 18 November 2025.
\url{https://arxiv.org/pdf/2310.20678v3}.
See especially Section 2, Proposition 2.11, and the stated arithmetic
nonvanishing hypotheses.

\bibitem{serban}
V. Serban, \emph{An infinitesimal \(p\)-adic multiplicative Manin--Mumford
conjecture}, Journal de Th\'eorie des Nombres de Bordeaux
30 (2018), no.~2, 393--408.
\url{https://www.numdam.org/item/JTNB_2018__30_2_393_0.pdf}.

\bibitem{gun-hazra-sahu}
S. Gun, S. Hazra and D. Sahu, \emph{On holomorphy and non-vanishing of
Artin \(L\)-functions}, Monatshefte f\"ur Mathematik 206 (2025), 853--883.
\url{https://doi.org/10.1007/s00605-025-02060-7}.

\bibitem{ltz2024}
R. J. Lemke Oliver, J. Thorner and A. Zaman,
\emph{An approximate form of Artin's holomorphy conjecture and
non-vanishing of Artin \(L\)-functions},
Inventiones Mathematicae 235 (2024), 893--971.
\url{https://doi.org/10.1007/s00222-023-01232-2}.
Author manuscript:
\url{https://lemkeoliver.github.io/papers/27-ApproximateArtin.pdf}.

\bibitem{bach-sorenson}
E. Bach and J. Sorenson, \emph{Explicit bounds for primes in residue
classes}, Mathematics of Computation 65 (1996), 1717--1735.
\url{https://doi.org/10.1090/S0025-5718-96-00763-6}.
Theorem 5.1 and Table 3 supply the explicit split-prime bounds used here.

\bibitem{ekn}
A.-S. Elsenhans, J. Kl\"uners and F. Nicolae,
\emph{Imaginary quadratic number fields with class groups of small
exponent}, arXiv:1803.02056.
\url{https://arxiv.org/pdf/1803.02056}.
Theorems 1 and 3 and Tables 1 and 2 supply the census and global range.

\bibitem{rodgers-tao}
B. Rodgers and T. Tao, \emph{The de Bruijn--Newman constant is
non-negative}, arXiv:1801.05914.
\url{https://arxiv.org/abs/1801.05914}.

\end{thebibliography}


\endgroup
\end{document}

